% Auto-typeset from the previously authored formula-by-formula study manuscript.
% Based on the thesis main branch f80e5ac903dd, 2026-10-31.

\chapter{公式与符号导航}

这是索引，不代替六章的逐式讲解。查到某个公式不理解时，直接打开对应的精读文件，再回原论文源码核对。\par

\section{一、章节与源码对应}

\begin{description}
\item[论文] 源码； 逐式精读
\item[第一章：绪论] ../hmb-nzbc-thesis/section/01-introduction.tex； \textbf{01-绪论逐句解析.md}
\item[第二章：模型与谱分析] ../hmb-nzbc-thesis/section/02-model.tex； \textbf{02-模型与非零背景谱分析.md}
\item[第三章：Jost 与散射] ../hmb-nzbc-thesis/section/03-scattering.tex； \textbf{03-Jost解与直接散射.md}
\item[第四章：矩阵 RHP] ../hmb-nzbc-thesis/section/04-rhp.tex； \textbf{04-矩阵RHP与有限维系统.md}
\item[第五章：一轨道直接求解] ../hmb-nzbc-thesis/section/05-solitons.tex； \textbf{05-一轨道解与逆重构.md}
\item[第六章：性质与共振] ../hmb-nzbc-thesis/section/06-conclusion.tex； \textbf{06-传播性质与共振正则性.md}
\end{description}

\section{二、重要数学对象与依赖}

\begin{description}
\item[对象] 作用； 首次重要位置
\item[q、p、\ensuremath{\eta}] 光场、介质极化、反转； eq:HMB（第一章）
\item[Q>0] 光场背景模长； eq:targetboundary
\item[\ensuremath{\eta}\textsubscript{0}、a、\ensuremath{\Omega}] 介质端态、共振参数、载波频率； eq:background（第二章）
\item[A(q)、R、\ensuremath{\sigma}\textsubscript{3}] 空间势、介质矩阵、对角约化矩阵； eq:lax 前
\item[U、V] 两部分 Lax 矩阵； eq:lax
\item[G、u、v] 去相位规范及新场变量； eq:gauge
\item[F(\ensuremath{\lambda})] 时间背景相位函数，含 \ensuremath{\lambda}=a 共振； eq:F
\item[\ensuremath{\rho}(\ensuremath{\lambda})] 背景谱根，\ensuremath{\rho}\textsuperscript{2}=\ensuremath{\lambda}\textsuperscript{2}+Q\textsuperscript{2}； eq:uniform
\item[w] 双叶曲面单值化坐标； eq:uniform
\item[h(w)] 换叶映射 \ensuremath{-}Q\textsuperscript{2}/w； eq:sheet
\item[D\textsubscript{+}、D\textsubscript{-}] 按 Im\ensuremath{\rho} 正负划分的解析区域； eq:domains
\item[E、\ensuremath{\gamma}] 背景特征向量矩阵及其行列式； eq:E、eq:diagE
\item[\ensuremath{\theta}] 背景同时 Lax 相位 \ensuremath{\rho}[x\ensuremath{-}Fz]； eq:theta
\item[\ensuremath{\chi}\ensuremath{\pm}、\ensuremath{\mu}\ensuremath{\pm}、\ensuremath{\Psi}\ensuremath{\pm}] 空间、去振荡、同时归一化 Jost 解； 第三章
\item[S\_sp、S] 空间/同时归一化散射矩阵； eq:Sdef
\item[s\textsubscript{1}\textsubscript{1}、s\textsubscript{2}\textsubscript{2}] 解析透射系数，零点给离散谱； eq:wronskian
\item[r、\textbackslash{}widetilde r] 反射系数； eq:reflection
\item[\ensuremath{\zeta}\_j、b\_j、c\_j、C\_j] 离散点、列相关常数、范数常数与带时空指数的系数； 第三章末
\item[M(w)] RHP 待求矩阵，由解析列拼接； eq:Mdef
\item[\textbackslash{}mathcal A、\textbackslash{}mathcal B] 有限维留数系统的系数矩阵； eq:finite-schur
\item[b、\ensuremath{\kappa}、W、h\textsubscript{0}、t] 第五章一轨道的计算参数； eq:imaginarypoint、eq:orbit-phase
\item[P、N、H] 直接解得到的投影缩写、分子矩阵、相位矩阵； eq:algebraic-P、eq:Nfactor
\item[\textbackslash{}mathcal U、\textbackslash{}mathcal V] 逆构造中 \ensuremath{\Psi} 的对数导数； eq:inverse-log-lax
\item[A\_b、B\_b] 局域波速度与内部频率； eq:velocity、eq:frequency
\end{description}

\section{三、逐式标签快速检索}

\subsection{第一章}

\begin{itemize}[leftmargin=2em]
\item eq:HMB：原 HMB 三方程、全部变量。
\item eq:targetboundary：左右同相非零光场及介质端态。
\end{itemize}

\subsection{第二章}

\begin{itemize}[leftmargin=2em]
\item prop:bg / prop:bloch：平面波与 Bloch 长度守恒。
\item eq:background / eq:bloch / eq:C：背景、x 守恒、长度常数。
\item eq:lax / eq:V32 / eq:V1 / eq:V0：空间矩阵与时间三次多项式全部系数。
\item prop:zc / eq:zc / eq:antihermitian / eq:rationalzc：零曲率、反 Hermitian 对称、介质方程。
\item eq:gauge / eq:gaugelax / eq:Rbg / eq:F：规范变换与背景时间相位。
\item eq:uniform / eq:lambda-spectrum / eq:Sigma / eq:domains / eq:sheet：单值化与谱区域。
\item eq:E / eq:diagE / eq:theta / eq:resonance-w：背景特征向量、相位、共振两原像。
\end{itemize}

\subsection{第三章}

\begin{itemize}[leftmargin=2em]
\item thm:jost / eq:volterra / eq:columnkernel：Jost 构造与解析列。
\item eq:Sdef / prop:scattering / eq:detjost / eq:wronskian：散射矩阵与透射、反射系数。
\item eq:conjugation / eq:musheet / eq:sconj / eq:Esheet / eq:reflection：约化与散射数据关系。
\item prop:quartet / eq:quartet：四点离散谱轨道。
\item prop:evolution / eq:timeS / eq:timer：传播时空演化。
\item prop:resonance：共振点零反射的条件性必要结果。
\item prop:highenergy / eq:mu-highenergy / eq:transmission-highenergy / eq:smallmu / eq:n-volterra-system：两个无穷远渐近及重构基础。
\end{itemize}

\subsection{第四章}

\begin{itemize}[leftmargin=2em]
\item thm:RHP / eq:Mdef / eq:jump：解析列拼接与跳跃。
\item eq:detM / eq:Minfty / eq:Mzero / eq:Msym：矩阵约化与双归一化。
\item eq:resplus / eq:resminus / eq:cpair：离散谱留数与配对。
\item prop:finite / eq:partialfraction / eq:finite1 / eq:finite2 / eq:finite-schur：无反射有限维系统。
\item prop:theta / eq:trace / eq:theta-condition：迹公式与同相端态。
\item prop:reconstruction / eq:reconq / eq:RfromM / eq:Mspace / eq:qfinite：场的重构。
\end{itemize}

\subsection{第五章}

\begin{itemize}[leftmargin=2em]
\item eq:imaginarypoint / eq:onezeros / eq:onec / eq:orbit-phase：一轨道离散数据。
\item lem:direct-solve / eq:one-AB / eq:one-schur-det / eq:direct-vu1 / eq:direct-vu2：显式消元、全局非零分母。
\item eq:algebraic-P / eq:Nfactor / thm:oneRHP / eq:explicitM：直接求解得到的 RHP 矩阵。
\item thm:direct-HMB / eq:finalfields / eq:etacomponent / eq:pcomponent / eq:direct-medium：三场重构。
\item eq:Pends / eq:direct-M-ends / eq:inverse-log-lax / eq:inverse-V-form：同相端态与 Lax 逆重建。
\item eq:hyperbolic-projector：一轨道波形的双曲/三角表示。
\item eq:N-cramer / eq:exactparams：一般有限系统的条件性表示与精确参数例。
\item prop:darboux / eq:explicit-y / eq:projector / eq:darbouxfields / eq:Nrecursion / eq:Nmedium：与已有 Darboux 的关系；不是直接 RHP 求解的起点。
\end{itemize}

\subsection{第六章}

\begin{itemize}[leftmargin=2em]
\item eq:velocity / eq:frequency：F(ib) 的实、虚部。
\item prop:wave-properties / eq:amplitude-bound / eq:medium-bound / eq:internal-period：波幅、传播、周期。
\item prop:oscillatory-regularity / eq:oscillatory-derivatives：奇异振荡函数的 C\textasciicircum{}k 延拓判据。
\item prop:resonant-regularity / eq:local-resonant-phase：HMB 共振反射数据的局部消失阶条件。
\end{itemize}

\section{四、重要的依赖关系}

\textbf{主线 A：构造 Jost 解}\par

eq:background \ensuremath{\to} eq:gauge \ensuremath{\to} eq:F \ensuremath{\to} eq:uniform \ensuremath{\to} eq:E \ensuremath{\to} eq:theta \ensuremath{\to} eq:volterra \ensuremath{\to} thm:jost。\par

\textbf{主线 B：构造 RHP}\par

thm:jost \ensuremath{\to} eq:Sdef \ensuremath{\to} eq:wronskian \ensuremath{\to} eq:conjugation 与 eq:musheet \ensuremath{\to} prop:quartet \ensuremath{\to} eq:Mdef \ensuremath{\to} eq:jump \ensuremath{\to} eq:resplus、eq:resminus、eq:cpair。\par

\textbf{主线 C：直接求一轨道}\par

eq:partialfraction \ensuremath{\to} eq:finite1、eq:finite2 \ensuremath{\to} eq:finite-schur \ensuremath{\to} eq:onezeros、eq:onec \ensuremath{\to} eq:one-AB \ensuremath{\to} eq:one-schur-det \ensuremath{\to} eq:direct-vu1、eq:direct-vu2 \ensuremath{\to} eq:explicitM。\par

\textbf{主线 D：恢复 HMB}\par

eq:explicitM \ensuremath{\to} eq:reconq 与 eq:RfromM \ensuremath{\to} eq:finalfields \ensuremath{\to} eq:inverse-log-lax \ensuremath{\to} eq:inverse-V-form \ensuremath{\to} eq:zc \ensuremath{\to} eq:HMB。\par

\textbf{主线 E：共振正则性}\par

eq:Rbg \ensuremath{\to} prop:resonance \ensuremath{\to} eq:timer \ensuremath{\to} prop:oscillatory-regularity \ensuremath{\to} prop:resonant-regularity。\par

\section{五、学习时应该特别回避的概念混淆}

\begin{enumerate}[leftmargin=2em]
\item \textbf{空间谱函数}与\textbf{传播演化矩阵}属于两个 Lax 方向；共振点 a 使后者奇异，但前者在 a 可以正则。
\item \textbf{整个 Jost 矩阵有定义}与\textbf{两列在某复域解析}是不同表述。
\item \textbf{r\ensuremath{\equiv}0} 与 \textbf{r(a)=0} 不同；前者是整条连续谱上的无反射问题，后者仅是一个点的必要相容条件。
\item \textbf{一轨道}含完整四点，m=2 是 D\textsubscript{+} 的点数，并不等于两个完全自由的孤子。
\item \textbf{Schur 系统有解} 与 \textbf{RHP 及其逆 Lax 全部条件验证成功}不同。
\item \textbf{与 Darboux 解代数一致} 不等于 \textbf{本论文从 Darboux 起步}；当前第五章已经先求直接 RHP，后做 Darboux 比较。
\item \textbf{局部 C\textasciicircum{}k 判据} 不等于 \textbf{一般含反射数据的全局 IST 可解性}。
\end{enumerate}

\section{六、检索与版本说明}

本导航使用原论文源码中的 LaTeX 标签作为定位方式；不依赖 PDF 页码，因为排版修订会改变页码。对应基准为 2026-10-31 提交 f80e5ac903dd58ee150fb523a14297e97e21c7e3。若之后论文新增或删除公式，应按标签重新核对并更新本目录。\par

本目录的数学解释根据论文给出的条件及证明展开，\textbf{不是对所有定理的外部独立审稿}。遇到分支点可去性、全局唯一性、一般反射的局部空间等高级结论时，应结合源码重新逐行验证其所用假设。\par

